\chapter{空間と座標系の関係}
    空間と座標系の関係は次の通りである。必ずしも包含関係ではないことに留意すべし。
    $N$ は自然数である。
    \begin{enumerate}
        \item $\realNumbers^N\coloneq N$ 次元実ベクトル空間
        \item $\mathbb{E}^N（N \text{次元 Euclid 空間}）\coloneq\realNumbers^N$ に標準内積が付加されたもの。例えば物理学では，ここが物理現象の舞台。
        \item 座標系：$\EuclideanSpace^N$ と $\realNumbers^N$ 間の（特異点を除いて）可逆な写像
        \item Euclid 空間上の，座標が陽に現れない関係式が，ある座標系上で成り立つなら，他の座標系上でも成り立つ。
        \item スカラー場：まず $\EuclideanSpace^N$ から $\realNumbers$ への写像として定義される。しかし実際の計算に於いては，座標系を定めて $\realNumbers^3\to\EuclideanSpace^3\to\realNumbers$ の合成写像が使われるのが普通である。
        \item ベクトル場：まず $\EuclideanSpace^N$ から $\EuclideanSpace^N$ への写像として定義される。スカラー場と同様の留意事項あり。
        \item 勾配，発散は $\EuclideanSpace^N$ 上で定義される。座標系がそれを $\realNumbers^N$ に写す。
        \item 外積，回転は $\EuclideanSpace^3$ 上で定義される。座標系がそれを $\realNumbers^3$ に写す。
    \end{enumerate}